
@@M0@@-module localement libre de rang @@M1@@, il suffit de montrer, compte tenu de 5.2.2 b), que
@@M2@@
Comme @@M3@@ est bornée, l'objet final @@M4@@ de @@M5@@ se décompose en une somme directe \emph{finie} disjointe @@M6@@ de sous-objets de telle façon que @@M7@@ est constante sur chaque sous-topos @@M8@@ @@M9@@. Comme le foncteur @@M10@@ @@M11@@ transforme sommes directes \emph{finies} en produits directs finis, on peut supposer que @@M12@@ est constant sur @@M13@@, soit @@M14@@. Posant @@M15@@, on trouve (cf. Exp V) :
@@M16@@
ce qui démontre (5.3.3), le premier membre étant égal à
@@M17@@

Il nous reste à démontrer le

\textbf{Lemme 5.3.4.} \uline{Pour tout topos annelé} @@M18@@ \uline{et pour tout} @@M19@@\uline{, on a}
@@M20@@

\emph{Démonstration} : Soit @@M21@@ l'homomorphisme d'augmentation. Alors (6.2.8) et (6.5.3) ci-dessous donnent par linéarité la formule
@@M22@@
